
\section{Fibonacci bounds and logarithmic cutoff estimates}
\label{app:fibonacci_bounds}

This appendix records standard Fibonacci estimates used throughout the quantitative scaling arguments.

\subsection{Binet formula and elementary bounds}
Let $\varphi=\frac{1+\sqrt{5}}{2}$ and $\psi=\frac{1-\sqrt{5}}{2}=-\varphi^{-1}$.
Then the Fibonacci numbers $(F_n)$ with $F_1=F_2=1$ satisfy the Binet formula
\[
F_n=\frac{\varphi^n-\psi^n}{\sqrt{5}}
\qquad (n\ge 1),
\]
and in particular $F_n=\Theta(\varphi^n)$.

\begin{lemma}[Two-sided bounds]
\label{lem:fib_two_sided_bounds_app}
For all $n\ge 2$ one has
\[
\varphi^{\,n-2}\le F_n\le \varphi^{\,n-1}.
\]
\end{lemma}

\begin{proof}
For $n=2,3$ this is immediate.
Assume $n\ge 4$ and that the bounds hold for $n-1$ and $n-2$.
Using $F_n=F_{n-1}+F_{n-2}$ and $\varphi^2=\varphi+1$ gives
\[
F_n\le \varphi^{n-2}+\varphi^{n-3}=\varphi^{n-3}(\varphi+1)=\varphi^{n-1},
\]
and similarly
\[
F_n\ge \varphi^{n-3}+\varphi^{n-4}=\varphi^{n-4}(\varphi+1)=\varphi^{n-2}.
\]
\end{proof}

\subsection{Admissible density and the cutoff index}
The admissible language size is $|X_m|=F_{m+2}$, hence the admissible density satisfies
\[
\frac{|X_m|}{2^m}=\frac{F_{m+2}}{2^m}=\Theta\!\Bigl(\Bigl(\frac{\varphi}{2}\Bigr)^m\Bigr),
\]
and the mean folding degeneracy is $\overline{d}_m=2^m/F_{m+2}=\Theta((2/\varphi)^m)$.

The cutoff index $K(m)$ defined by $F_{K(m)+1}\le 2^m-1<F_{K(m)+2}$ satisfies $K(m)=m\log_\varphi 2+O(1)$, as shown in Proposition~\ref{prop:Km_bounds}.


